% =====================================================================
% Owner: 5   Budget: ~2.0 pages
% Every number below was re-checked on 2026-11-09 against the final evidence
% package (DSA4213_final_20261028: official harness report.json per run,
% run_result.json, labels.json, prompt_candidates.json) and the Role 5
% analysis pages (tmp/claude-formal-charts-review-20261028).
% Per-run first-view data: data/first_reference_view.csv (from the Role 5 analysis page 06_exploration).
% =====================================================================
\section*{5. Evaluation and Results}

\paragraph{Protocol.}
All conditions ran on the same 40 held-out tasks with the same model, instructions, Docker images
and budgets; only the localization context differs (Table~\ref{tab:setup}). The proposal's primary
comparison is \condA{} vs \condD; \condB{} and \condC{} isolate the lexical and structural components.
\emph{Repair:} a run is resolved when the official harness~\cite{jimenez2024swebench,swebenchharness}
passes every \texttt{FAIL\_TO\_PASS} (F2P, target) and \texttt{PASS\_TO\_PASS} (P2P, existing) test.
Of the 224 runs, 216 were evaluated on the submitted patch; two (sympy-12489 \condB, \condC) were
re-evaluated on the original submission because the recorded patch file differed from it; and six
ended without a submission (e.g.\ repeated output-format errors) and were evaluated on the patch
extracted from their final state. Only one of the six passed (sympy-13757 \condB); counting all six as
unresolved changes \condB{} from 23 to 22 and nothing else.
\emph{Localization:} the usual file-level metrics~\cite{zhou2012buglocator,chen2025locagent}
(Hit@1/5/20, Recall@5, MRR) of the initial ranking, with the reference patch's Python source files as
targets (tests and documentation excluded; \condA{} has no ranking), plus the first model call whose
shell command reads a target file (\texttt{cat}, \texttt{sed}, \texttt{grep}, \dots; an approximation).
\emph{Efficiency:} model calls, tokens (including format-error retries), API cost and agent seconds;
ranking time is measured separately (Section~6).
\emph{Variability:} one main-round run per task and condition; 8 tasks, drawn before any result was
seen, ran two further rounds in every condition (64 runs). Paired differences use an exact McNemar test.

\begin{table}[h]
\centering\small
\caption{Main results on the 40 held-out tasks (one run per task and condition).
F2P = target tests; P2P = existing tests.}
\label{tab:main}
\begin{tabular}{@{}lcccc@{}}
\toprule
 & \condA{} original & \condB{} +text & \condC{} +structure & \condD{} +failures \\
\midrule
Resolved (/40)                          & 20 & 23 & \textbf{25} & 21 \\
F2P pass, P2P broken                    & 5  & 2  & 1  & 3  \\
F2P fail, P2P intact                    & 12 & 12 & 11 & 13 \\
F2P fail, P2P broken                    & 3  & 3  & 3  & 3  \\
\midrule
Reference file viewed at call 1 (/40)   & 6  & 29 & \textbf{30} & 29 \\
Model calls (total)                     & 367 & 357 & 343 & 367 \\
Total tokens (M)                        & 3.18 & 3.76 & 3.69 & 3.64 \\
Tokens per resolved task (k)            & 159 & 163 & \textbf{148} & 173 \\
Cost (\$) / cost per resolved task (\$) & 0.86 / 0.043 & 1.22$^*$ / 0.053 & 1.05 / 0.042 & 0.95 / 0.045 \\
Median agent seconds per run            & 19 & 17 & 16 & 19 \\
\bottomrule
\end{tabular}\\[2pt]
{\footnotesize $^*$One \condB{} run (sphinx-9229) has incomplete cost records, so \condB's cost is a lower bound; token counts are complete.}
\end{table}

\paragraph{RQ1: repair.}
No hint reduced repair: \condC{} resolved the most tasks (25/40), \condA{} the fewest (20/40), \condB{}
and \condD{} 23 and 21. The ordering rests on 7 \emph{swing} tasks, since 19 tasks were resolved in every
condition and 14 in none (Appendix~A.3). \condC{} resolved six tasks that \condA{} did not and \condA{} one
that \condC{} did not (exact McNemar $p=0.13$); the pre-specified \condA{} vs \condD{} split is 1 vs 2
($p=1.0$). Two observations place these gaps within run-to-run noise. On the 30 tasks where \condD's
feedback never fired, its prompt is identical to \condC's, yet 3 outcomes differ. In the stability
subset the order becomes \condD{} 17, \condB{} 15, \condC{} 15, \condA{} 14 of 24, and every condition
flips at least one task across rounds (Appendix~A.5). We therefore claim that the hints did not hurt
repair, not that any condition repairs better.

The conditions differ less in fixing the target than in what else they break. Target tests pass in
25, 25, 26 and 24 runs for \condA--\condD, but \condA{} breaks existing tests in 5 of them and \condC{} in 1.
In pytest-10081 both edit the same condition in \texttt{src/\_pytest/unittest.py}: \condA{} replaces the
skip check with one on the parent class and breaks 1 of 63 existing tests, while \condC{} keeps the
original check, adds the parent's and passes all 63. Passing the target tests is thus not the same as
a correct patch, a known risk in automated program repair~\cite{smith2015cure,yu2019overfitting}.
Whether structure causes fewer regressions is untested; we report it only as a lead.

\paragraph{RQ2: exploration cost.}
Hints bring a reference file into view about one model call earlier: at the first call the agent had
read one in 29--30 of 40 tasks with a hint and in 6 without, and by the second call \condA{} has caught
up (31 vs 31--34). The hint thus replaces the initial search, not later work. Hinted runs used more
tokens in total (3.64--3.76M vs 3.18M), because the ${\sim}550$-token hint stays in the history and is
re-sent with every call (5.5\% of input tokens). Per resolved task \condC{} was the cheapest (148k vs
159k for \condA), and on the six tasks only \condC{} solved it used 330k tokens to \condA's 454k. The
saving comes from harder tasks: on the 19 tasks every condition solved, the hint is pure overhead
(mean 55k vs 40k tokens per run for \condC{} vs \condA). Against the success criterion of Section~1,
\condC{} gives comparable repair at a lower cost \emph{per fix}, but not at a lower total cost.

\subsection*{5.1 Comparisons and Ablations}

\begin{table}[h]
\centering\small
\caption{Initial ranking quality on the 40 held-out tasks (counts of tasks unless stated); \condD's
initial ranking is identical to \condC's. ``In hint'': a target is among the 20 source files shown
(tests are filtered out first, so this can exceed Hit@20); ``Viewed'': the agent read a target, among
the Hit@20 tasks.}
\label{tab:loc}
\begin{tabular}{@{}lccccccc@{}}
\toprule
 & Hit@1 & Hit@5 & Recall@5 & MRR & Hit@20 & In hint & Viewed \\
\midrule
\condB{} (BM25 $L$)            & 9 & 16 & 38.8\% & 0.320 & 27 & 30 & 27/27 \\
\condC{} ($\mathrm{RRF}(L,S)$) & 9 & 16 & 38.8\% & 0.322 & 28 & 30 & 28/28 \\
\bottomrule
\end{tabular}
\end{table}

\paragraph{\condB{} vs \condC: what structure adds.}
\condC{} keeps \condB's top~10 (Section~4.2), so Hit@1 and Hit@5 are identical and MRR differs by 0.002.
Structure acts only on ranks 11--20, adding one task to Hit@20 (28 vs 27; Django $+2$, SymPy $-1$). The
clearest gain is django-16950: its reference file moves from rank 35 to 19 and enters \condC's hint;
both conditions resolved it, \condC{} with 62k tokens and \condB{} with 127k. On repair, \condC{} resolved
two tasks that \condB{} did not and never the reverse, too few to attribute to structure.

\paragraph{\condC{} vs \condD: what failure feedback adds.}
\condD{} resolved four tasks fewer than \condC{} and none that \condC{} missed (django-14311,
django-16950, sphinx-7757, pytest-7521). Feedback reached the agent in only one of these runs, and only
as a diagnostic note (django-14311), so the gap is mostly noise rather than an effect of feedback.
Overall, feedback was delivered in 10 of 40 \condD{} runs and re-ranked files in 5; Section~6 explains
why it rarely changed the outcome.
